\documentclass[11pt,a4paper]{article}
\usepackage[utf8]{inputenc}
\usepackage[T1]{fontenc}
\usepackage{lmodern}
\usepackage{amsmath,amssymb}
\usepackage{booktabs}
\usepackage{xcolor}
\usepackage{geometry}
\usepackage[colorlinks=true,linkcolor=blue!50!black,urlcolor=blue!50!black]{hyperref}
\geometry{a4paper,left=2.2cm,right=2.2cm,top=2.2cm,bottom=2.4cm}
\setlength{\parindent}{0pt}
\setlength{\parskip}{0.7ex}
\numberwithin{equation}{section}
\newcommand{\code}[1]{\texttt{\small #1}}
\newcommand{\dg}{^{\dagger}}
\newcommand{\tr}{\operatorname{tr}}

\title{The CUDA kernels of RutticeQCD, one by one}
\author{Wei-Lun Chen}
\date{\today\quad(internal note)}

\begin{document}
\maketitle

\noindent For every kernel: what it computes, written out line by line, how its threads are laid
out, and what in memory decides its speed. File names are under \code{cuda\_bridge/include/}.
The communication kernels are in the note on several GPUs.

\tableofcontents

%======================================================================
\section{How the fields lie in memory}

\subsection{Blocks of 32 sites}

A field has $N_c$ reals per site. Sites are grouped by 32, and inside a group the same component
of the 32 sites is consecutive:
\begin{equation}
a(c,n) = \big(\lfloor n/32\rfloor\cdot N_c + c\big)\cdot 32 + (n \bmod 32).
\label{eq:aosoa}
\end{equation}
Why: a warp is 32 threads. If thread $j$ handles site $32k+j$ and reads component $c$,
\begin{equation}
a(c,32k+j) = (kN_c+c)\cdot32 + j,\qquad j=0,\dots,31,
\end{equation}
32 consecutive numbers: one memory transaction. This is the only rule behind every layout below.

\subsection{Spinors with a fifth dimension}

A domain-wall spinor has $L_s$ spinors of 24 reals per site. Component (colour $k$, spin $a$,
real/imag $r$) at fifth coordinate $s$ is component
\begin{equation}
c = 24s + 2(k + 3a) + r
\end{equation}
of the site, $N_c=24L_s$. A field on one parity numbers its sites by $\lfloor x/2\rfloor$:
\begin{equation}
n_{\text{half}} = \lfloor x/2\rfloor + \tfrac{L_x}{2}\big(y+L_y(z+L_zt)\big),\qquad
x = 2\lfloor x/2\rfloor + \big((p+y+z+t)\bmod 2\big),
\end{equation}
the second line recovering $x$ from the parity $p$. Several vectors of a field lie one after
another, each padded to a whole number of blocks.

\subsection{Links: two rows are enough}

A link $U\in SU(3)$ has rows $r_0,r_1,r_2$. Unitarity and $\det U=1$ give
\begin{align}
r_2 &= \overline{r_0\times r_1},\\
(r_2)_c &= \overline{(r_0)_a(r_1)_b - (r_0)_b(r_1)_a},\qquad a=c+1,\ b=c+2\pmod 3 .
\end{align}
So the operators keep 12 reals per link instead of 18 when the links are in $SU(3)$, and rebuild
row 2 in registers (\code{load\_link} in \code{hops.cuh}). A hop reads mostly links, so this
saves about a third of what it reads.

%======================================================================
\section{The domain-wall operator}

\subsection{What the operator is}

The Möbius operator on $\psi_s$, $s=0,\dots,L_s-1$, is
\begin{equation}
D\psi_s = (b\,D_W+1)\psi_s + (c\,D_W-1)\big(P_-\psi_{s+1}+P_+\psi_{s-1}\big),
\end{equation}
with $D_W = (4-M_5) - \tfrac12\tilde H$ and $\tilde H$ the hops without $\kappa$. Define the
fifth-dimension hop
\begin{equation}
(Y\psi)_s = -P_-\psi_{s+1} - P_+\psi_{s-1}\qquad(+m\text{ in place of }-1\text{ across the walls}).
\end{equation}
Then $P_-\psi_{s+1}+P_+\psi_{s-1} = -(Y\psi)_s$, and
\begin{align}
D\psi &= (bD_W+1)\psi - (cD_W-1)Y\psi\\
&= \big[b(4-M_5)+1\big]\psi + \big[1-c(4-M_5)\big]Y\psi - \tfrac12\tilde H\,(b\psi - cY\psi)\\
&= \underbrace{(d_+ + d_-Y)\psi}_{C\psi} \;-\; \tfrac12\tilde H\,(b - cY)\psi,
\end{align}
with $d_+ = 1+b(4-M_5)$ and $d_- = 1-c(4-M_5)$. $C$ acts within a four-dimensional site; the
hops act on $w=(b-cY)\psi$. Every kernel below is a piece of this line.

\subsection{\code{fifth\_term}: $a\psi_s + y(Y\psi)_s$}

\begin{align}
w_s &= a\,\psi_s + y\,(Y\psi)_s\\
&= a\,\psi_s - y\,P_-\psi_{s+1} - y\,P_+\psi_{s-1}\qquad (0<s<L_s-1),\\
w_0 &= a\,\psi_0 - y\,P_-\psi_1 + y\,m\,P_+\psi_{L_s-1},\\
w_{L_s-1} &= a\,\psi_{L_s-1} + y\,m\,P_-\psi_0 - y\,P_+\psi_{L_s-2}.
\end{align}
For $Y\dg$ the projectors swap: $P_+$ from above, $P_-$ from below.

\subsection{\code{hop\_kernel}: $\text{out} = [\text{from}-]\,H(b-cY)\,\text{in}$}

At output site $x$ and fifth coordinate $s$:
\begin{equation}
\text{out}_s(x) = -\kappa\sum_\mu\Big[(1-s_\pm\gamma_\mu)U_\mu(x)\,w_s(x+\hat\mu) + (1+s_\pm\gamma_\mu)U_\mu(x-\hat\mu)\dg\,w_s(x-\hat\mu)\Big],
\end{equation}
with $w_s$ from \code{fifth\_term} computed on the fly at the neighbour ($a=b$, $y=-c$) and
$s_\pm=\pm1$ for $H$, $H\dg$. A leg is done in three steps (\code{leg\_half}, link,
\code{add\_leg}). For $\mu=0$ (Sec.~1.1 of the other note):
\begin{align}
h_0 &= w_0 + is_\pm w_3, & h_1 &= w_1 + is_\pm w_2,\\
t_a &= U\,h_a\quad(a=0,1), \\
\text{out}_0 &\mathrel{+}= -\kappa\,t_0, & \text{out}_1 &\mathrel{+}= -\kappa\,t_1,\\
\text{out}_3 &\mathrel{+}= -\kappa\,(-is_\pm)\,t_0, & \text{out}_2 &\mathrel{+}= -\kappa\,(-is_\pm)\,t_1 .
\end{align}
Two $SU(3)$ products per leg instead of four.

\textbf{Threads.} A thread for each (site, fibre), a fibre being one $s$ of one vector. The
threads of a block are 2D: \code{threadIdx.x} the site (consecutive sites: coalesced by
Eq.~\eqref{eq:aosoa}), \code{threadIdx.y} the fibre. All fibres of a site read the same eight
links, which then come from the cache after the first. At most 256 threads a block: a thread may
use 255 registers and a block has 65536.

\textbf{Double precision.} The two legs of a direction are done one after the other, so only one
spinor and one link are live: 168 registers, three blocks per multiprocessor. The sums go into
\code{out} in the same order either way.

\subsection{\code{mdwf\_dslash\_1pass\_kernel}: all of $D$ in one pass}

\begin{equation}
\text{out}_s(x) = \underbrace{d_+\psi_s(x) + d_-(Y\psi)_s(x)}_{\code{fifth\_term}(a=d_+,\,y=d_-)} \;+\; \text{hops}\big(\kappa=\tfrac12\big).
\end{equation}
One read of $\psi$ and the links, one write of out. Used for $D$ on the whole lattice.

\subsection{\code{mdwf\_fifth\_kernel}: $(a_1+y_1Y)x + (a_2+y_2Y)z$}

Two \code{fifth\_term}s added, for the parts of the even-odd operators that act within a site.

\subsection{The fifth dimension as a matrix: \code{fifth\_gemm}}

In the Dirac representation $\gamma_5$ swaps the upper and lower spin pairs:
\begin{equation}
\gamma_5=\begin{pmatrix}0&1\\1&0\end{pmatrix},\qquad
P_\pm = \tfrac12(1\pm\gamma_5) = \tfrac12\begin{pmatrix}1&\pm1\\\pm1&1\end{pmatrix}.
\end{equation}
Write a spinor at $s$ as (upper $u_s$, lower $l_s$). Any operator of the form
$\sum_{s'}(A_{ss'}P_+ + B_{ss'}P_-)$ gives
\begin{align}
\text{out}^u_s &= \sum_{s'}\Big[A_{ss'}\tfrac{u_{s'}+l_{s'}}{2} + B_{ss'}\tfrac{u_{s'}-l_{s'}}{2}\Big]
= \sum_{s'}\Big[\tfrac{A+B}{2}u_{s'} + \tfrac{A-B}{2}l_{s'}\Big]_{ss'},\\
\text{out}^l_s &= \sum_{s'}\Big[A_{ss'}\tfrac{u_{s'}+l_{s'}}{2} - B_{ss'}\tfrac{u_{s'}-l_{s'}}{2}\Big]
= \sum_{s'}\Big[\tfrac{A-B}{2}u_{s'} + \tfrac{A+B}{2}l_{s'}\Big]_{ss'}.
\end{align}
So the operator is one real $2L_s\times2L_s$ matrix $M$ acting on the rows $(u_0,l_0,u_1,l_1,\dots)$.
$Y$, $C$, $C^{-1}$, $b-cY$, Neff's $Q$, their adjoints and their products are all of this form.
Products are formed on the host, and one kernel applies the result. In a block of 32 sites the
rows of the field are $2L_s$ rows of $12\times32=384$ reals:
\begin{equation}
\text{out}\,(2L_s\times384) = M\,(2L_s\times2L_s)\ \text{in}\,(2L_s\times384).
\end{equation}

\textbf{Double precision on the tensor cores} (\code{fifth\_gemm\_mma\_kernel}). The instruction
\code{mma.m8n8k4.f64} computes $D(8\times8)\mathrel{+}=A(8\times4)\,B(4\times8)$; thread $t$ of a
warp holds $A[t/4][t\bmod4]$, $B[t\bmod4][t/4]$ and two entries of $D$. A warp takes a strip of 8
columns; for every 4 rows of the input it reads its $4\times8$ reals once and multiplies them
with the tiles of $M$.

\textbf{Why the XOR.} $M$ sits in shared memory, which has 32 banks of 4 bytes. A double occupies
two banks, so double number $k$ of a row of length $K_{16}$ (a multiple of 16) is in banks
\begin{equation}
\text{bank}(row,k) = 2(row\cdot K_{16}+k) \bmod 32 = 2k \bmod 32\qquad(2K_{16}\equiv0\bmod 32).
\end{equation}
A tile read takes rows $0..7$ and columns $k_0..k_0+3$: all eight rows hit the same banks, an
eight-way conflict. Store $(row,k)$ at column $k' = k \oplus 4(row\bmod4)$:
\begin{align}
row\bmod4=0:&\ k'\in\{k_0..k_0+3\}, & row\bmod4=1:&\ k'\in\{k_0..k_0+3\}\oplus4,\\
row\bmod4=2:&\ k'\in\{k_0..k_0+3\}\oplus8, & row\bmod4=3:&\ k'\in\{k_0..k_0+3\}\oplus12,
\end{align}
four disjoint groups of four doubles, $4\times4\times2=32$ banks. Rows $r$ and $r+4$ still share,
a two-way conflict. That is the least possible: 32 threads reading 8 bytes need 256 bytes, two
passes of 128.

\textbf{Single precision} (\code{fifth\_gemm\_columns\_kernel}): a thread per column of a block;
it reads the $2L_s$ reals of its column once and keeps the $8N_T$ results in registers. All
threads read the same entry of $M$ at the same time, a broadcast. \code{fifth\_gemm\_tf32\_kernel}
does the product on the tensor cores in TF32 (10 bits of mantissa, relative error $10^{-3}$), only
for operators whose result is corrected afterwards (smoothers, inner solves).

\subsection{\code{hop\_gemm\_kernel}: hops, then a matrix, in one pass}

$\text{out} = [\text{from}-]\,M\,H(b-cY)\,\text{in}$. The hops of a block of sites are written into
shared memory as the $2L_s$ rows of the block, then multiplied by $M$ as above. In double
precision the \code{mma} reads 4 rows $j..j+3$ and 8 columns of the tile. A row is $12\times$(sites
of the block) doubles; when that is a multiple of 16 doubles (128 bytes), rows $j$ and $j+1$ hit the
same banks. Storing column $k$ of row $j$ at
$k\oplus8(j\bmod2)$ shifts the odd rows by 8 doubles = 16 banks:
\begin{equation}
\text{bank}(j,k) = 2\big(k\oplus8(j\bmod2)\big)\bmod 32,
\end{equation}
so even and odd rows use the two halves of the banks; rows $j$ and $j+2$ share, two-way, again the
least for 32 doubles. This is the XOR seen in the code; the face buffers of several GPUs need none
(they are in global memory).

\subsection{Small kernels}

\begin{align}
\code{gamma5}:&\quad \text{out}_s = \gamma_5\,\text{in}_s,\ \text{or with reflection}\ \text{out}_s = \gamma_5\,\text{in}_{L_s-1-s}\ (\Gamma_5=\gamma_5R),\\
\code{walls} (\to4d):&\quad q = P_-\psi_{s_-} + P_+\psi_{s_+},\qquad (s_-,s_+)=(0,L_s-1)\ \text{for the propagator},\\
\code{walls} (\gets4d):&\quad \psi_{s_-} = P_-\eta,\ \psi_{s_+}=P_+\eta,\\
\code{parity\_copy}:&\quad \text{half}[n_{\text{half}}] \leftrightarrow \text{whole}[n].
\end{align}

%======================================================================
\section{Wilson and clover}

\subsection{\code{wilson\_interior\_kernel}}

$D\psi = \psi + H\psi$ with the hops of Sec.~2.3, $L_s=1$, $b=1$, $c=0$. The fibres of a site are
its vectors, which share the links.

\subsection{\code{clover\_term\_kernel}: $T\psi$}

$T = 1-\kappa c_{SW}\sigma\cdot F$ commutes with $\gamma_5$. In the Dirac representation the
eigenvectors of $\gamma_5$ are $(e_a\pm e_{a+2})/\sqrt2$, so in that basis $T$ is two Hermitian
$6\times6$ blocks (index colour $+3\times$spin):
\begin{align}
\psi^\pm_{k+3a} &= \psi_{a,k} \pm \psi_{a+2,k},\qquad a=0,1\quad(\text{the }1/\sqrt2\text{ left out}),\\
\phi^\pm &= T_\pm\,\psi^\pm,\\
(T\psi)_{a,k} &= \tfrac12\big(\phi^+_{k+3a} + \phi^-_{k+3a}\big),\qquad
(T\psi)_{a+2,k} = \tfrac12\big(\phi^+_{k+3a} - \phi^-_{k+3a}\big),
\end{align}
the two $1/\sqrt2$ making the $\tfrac12$. A Hermitian $6\times6$ block is stored packed: 6 real
diagonal entries and 15 complex above the diagonal,
\begin{equation}
6 + 2\cdot15 = 36\ \text{reals},\qquad 2\ \text{blocks} = 72\ \text{reals a site}.
\end{equation}
The same kernel applies $T^{-1}$ (its blocks made by \code{clover\_blocks\_kernel}, below).

\subsection{\code{clover\_dslash\_kernel}}

$D\psi = T\psi + H\psi$ in one pass.

\subsection{\code{clover\_blocks\_kernel}: $T_\pm$, $T_\pm^{-1}$ and $\ln|\det T_\pm|$ from the links}

One thread for each pair (site, sign $s=\pm1$); in double precision; the links extended one site
deep (Sec.~5). First the field strength of each plane $\mu<\nu$ from the four leaves at $x$,
\begin{align}
Q_{\mu\nu}(x) &= \sum_{\ell=1}^{4} W_\ell(x),\qquad
F_{\mu\nu}(x) = \tfrac18\big(Q_{\mu\nu}(x) - Q_{\mu\nu}(x)\dg\big),
\end{align}
$W_1$ the path $x\to x+\hat\mu\to x+\hat\mu+\hat\nu\to x+\hat\nu\to x$ and $W_2,W_3,W_4$ the other
three squares turning the same way (\code{leaves\_of}). Then
\begin{equation}
T = 1 - \kappa c_{SW}\sum_{\mu<\nu} G_{\mu\nu}\otimes F_{\mu\nu},\qquad G_{\mu\nu}=\gamma_\mu\gamma_\nu .
\end{equation}
$G_{\mu\nu}$ has one entry in each row: $G_{r,c(r)} = g(r)$. From the tables of $\gamma_\mu$,
\begin{equation}
c(r) = c_\nu\big(c_\mu(r)\big),\qquad g(r) = g_\mu(r)\,g_\nu\big(c_\mu(r)\big).
\end{equation}
The block of sign $s$ is $T$ between the vectors $v_{a}=(e_a+s\,e_{a+2})/\sqrt2$, $a=0,1$:
\begin{align}
(T_s)_{i+3a,\,j+3b} &= \delta_{ab}\delta_{ij} - \kappa c_{SW}\sum_{\mu<\nu} h^{s}_{ab}(\mu\nu)\,F_{\mu\nu,ij},\\
h^{s}_{ab} &= v_a\dg G\,v_b = \tfrac12\big(G_{ab} + s\,G_{a,b+2} + s\,G_{a+2,b} + G_{a+2,b+2}\big).
\end{align}
The kernel builds $h^s$ by going over the four rows $r$ of $G$: row $r$ adds $g(r)$ to
$h^s_{r\bmod2,\;c(r)\bmod2}$ with the factor $\tfrac12\,s^{[r\ge2]}\,s^{[c(r)\ge2]}$.

$T_s$ is packed (6 real diagonal entries, then the upper triangle) and inverted by Gauss–Jordan
with partial pivoting, as \code{numerics::invert\_complex}. Row operations that add a multiple of
one row to another do not change the determinant, and scaling the pivot row by $1/p$ divides it
by $p$, so
\begin{equation}
\ln|\det T_s| = \sum_{\text{columns}}\ln|p_{\text{col}}| .
\end{equation}
The change of basis to the $v_a$ is unitary, so $\ln|\det T(x)| = \ln|\det T_+| + \ln|\det T_-|$:
the even-odd action adds them over the sites of a parity (on the host, in site order). A pivot
below $6\,\epsilon\max|T_{rc}|$ marks the site singular, and \code{set\_links} returns an error.

Registers: $T_s$ and the inverse are $2\times72$ doubles, more than a thread keeps; the compiler
puts part of them in local memory (L1). The kernel runs once a step of the integrator, next to a
solve of hundreds of operator applications, so its speed does not matter.

\subsection{The force of the clover term: insertion, then leaves}

The clover term of $D$ is $-\kappa c_{SW}\sum_{\mu<\nu}F_{\mu\nu}\otimes G_{\mu\nu}$. Its part of
$-2\,\mathrm{Re}\,\zeta\dg\delta D\eta$, for $\delta U = XU$, is $\mathrm{Re}\tr(XM)$ with $M$
made in two steps.

\paragraph{\code{clover\_insertion\_kernel}.} For a pair $(\zeta,\eta)$ and a weight $w$, at every
site and plane,
\begin{align}
A_{\mu\nu,ij}(x) &= \sum_{s}(G_{\mu\nu}\eta)_{s,i}(x)\,\overline{\zeta_{s,j}(x)}
 = \sum_s g(s)\,\eta_{c(s),i}(x)\,\overline{\zeta_{s,j}(x)},\\
B_{\mu\nu}(x) &\mathrel{+}= \frac{w\,\kappa c_{SW}}{4}\big(A_{\mu\nu}(x) - A_{\mu\nu}(x)\dg\big).
\end{align}
One thread a site; the 24 reals of each spinor read once. $B$ of planes 0–3 sits in the four
directions of one field of links and that of planes 4–5 in the first two of another, so the halo
code of the links extends it.

\paragraph{\code{clover\_log\_det\_insertion\_kernel}.} The derivative of $-2\ln|\det T_o|$ (the
even-odd action) is the same insertion for the twelve pairs $(T^{-1}e_k, e_k)$ on the odd sites,
$e_k$ the unit vector of spin $\sigma$ and colour $c$. Only one component of $e_k$ is nonzero, so
the sum over $k$ collapses:
\begin{align}
A_{ij} &= \sum_{k}\sum_s g(s)\,(e_k)_{c(s),i}\,\overline{(T^{-1}e_k)_{s,j}}
 = \sum_s g(s)\,\overline{(T^{-1})_{(s,j),(c(s),i)}} .
\end{align}
The entries of the $12\times12$ $T^{-1}$ come from its packed blocks (\code{clover\_entry}):
\begin{equation}
(T^{-1})_{(a,i),(b,j)} = \sum_{s=\pm}\tfrac12\,s^{[a\ge2]}s^{[b\ge2]}\,(T_s^{-1})_{i+3(a\bmod2),\,j+3(b\bmod2)} .
\end{equation}

\paragraph{\code{clover\_leaves\_kernel}.} Once for all pairs: one thread a link $(x,\rho)$. Every
leaf of a plane containing $\rho$ that passes the link contributes. Cut the leaf at the link:
$W_1$ runs from the link to the leaf's corner, $B$ of that corner is inserted, $W_2$ runs back,
\begin{equation}
M_\rho(x) = \sum_{\text{leaves}} \pm\,W_1\,B_{\mu\nu}(\text{corner})\,W_2,
\end{equation}
$+$ if the leaf takes the link forward, $-$ if backward (\code{hmc/fermion\_action.rs},
\code{clover\_leaves}). Because $B$ is summed over the pairs first, a rational action with $N$
poles walks the leaves once instead of $N$ times; Bridge++ (\code{AForce\_F\_CloverTerm}) does the
eight terms $R(1)\ldots R(8)$ for every pair.

%======================================================================
\section{The fermion force: \code{hop\_force\_kernel}}

For $\delta U_\mu(x) = X\,U_\mu(x)$ the change of $-2\,\mathrm{Re}\,\zeta\dg\,\delta D\,\eta$ is
$\mathrm{Re}\,\tr(X M_\mu(x))$ with
\begin{equation}
M_\mu(x) = -2c\sum_s\sum_{\text{spin}}\big[U_\mu(x)(1-\gamma_\mu)\eta_s(x+\hat\mu)\big]\zeta_s(x)\dg
+ 2c\sum_s\sum_{\text{spin}}\big[\gamma_5\eta_s(x)\big]\big[U_\mu(x)(1-\gamma_\mu)\gamma_5\zeta_s(x+\hat\mu)\big]\dg .
\end{equation}
The kernel gets $U(1-\gamma_\mu)\eta(x+\hat\mu)$ from the forward leg of the hop kernel with
$\kappa=-1$ (so the leg has coefficient $+1$), and the second bracket the same way from
$\gamma_5\zeta$. Then the outer products, colour by colour:
\begin{equation}
(ab\dg)_{ij} = a_i\overline{b_j},\qquad
\mathrm{Re}=a^r_ib^r_j + a^i_ib^i_j,\quad \mathrm{Im}=a^i_ib^r_j - a^r_ib^i_j .
\end{equation}
A thread per site, a loop over $\mu$ and $s$; $M$ in the layout of the gauge fields (Sec.~5).

%======================================================================
\section{The gauge field}

\subsection{Layout}

Links and fields of matrices (staples, $\Sigma$, forces) are whole $3\times3$ complex matrices,
18 reals per direction, 72 per site, in blocks of 32 sites (Eq.~\eqref{eq:aosoa} with $N_c=72$).
All in double precision.

\subsection{\code{stout\_step\_kernel}: $V = e^{iQ}U$}

The staple sum and $Q$:
\begin{align}
C_\mu(x) &= \sum_{\nu\ne\mu}\rho_{\mu\nu}\Big[U_\nu(x)U_\mu(x+\hat\nu)U_\nu(x+\hat\mu)\dg + U_\nu(x-\hat\nu)\dg U_\mu(x-\hat\nu)U_\nu(x-\hat\nu+\hat\mu)\Big],\\
\Omega &= C\,U\dg,\\
iQ &= \tfrac12(\Omega-\Omega\dg) - \tfrac16\tr(\Omega-\Omega\dg)\cdot1\qquad(\text{traceless anti-Hermitian}).
\end{align}
The exponential by Cayley–Hamilton. $Q$ is Hermitian and traceless, so
\begin{equation}
Q^3 = c_1Q + c_0\cdot1,\qquad c_0 = \det Q = \tfrac13\tr Q^3,\qquad c_1 = \tfrac12\tr Q^2,
\end{equation}
and every power of $Q$ reduces to $1,Q,Q^2$:
\begin{equation}
e^{iQ} = f_0\cdot1 + f_1Q + f_2Q^2 .
\end{equation}
With the eigenvalues parametrised as
\begin{align}
\theta &= \arccos\frac{c_0}{c_0^{\max}},\qquad c_0^{\max} = 2\Big(\frac{c_1}{3}\Big)^{3/2},\\
u &= \sqrt{c_1/3}\,\cos\frac\theta3,\qquad w = \sqrt{c_1}\,\sin\frac\theta3,
\end{align}
Morningstar and Peardon find $f_j = h_j/(9u^2-w^2)$ with
\begin{align}
h_0 &= (u^2-w^2)e^{2iu} + e^{-iu}\big[8u^2\cos w + 2iu(3u^2+w^2)\,\xi_0(w)\big],\\
h_1 &= 2u\,e^{2iu} - e^{-iu}\big[2u\cos w - i(3u^2-w^2)\,\xi_0(w)\big],\\
h_2 &= e^{2iu} - e^{-iu}\big[\cos w + 3iu\,\xi_0(w)\big],\qquad \xi_0(w)=\frac{\sin w}{w}.
\end{align}
The kernel stores $iQ$ and $(iQ)^2$ and uses $Q=-i(iQ)$, $Q^2=-(iQ)^2$. When $|iQ|^2\le10^{-10}$ it
takes $e^{iQ}=1$. A thread for each (site, direction).

\subsection{\code{stout\_chain\_local\_kernel} and \code{stout\_chain\_staples\_kernel}}

The force on the thin links from $\Sigma'$, the force on the smeared ones (Morningstar–Peardon
eqs.~(70)--(75)):
\begin{align}
\Lambda &= \tfrac12(\Gamma+\Gamma\dg) - \tfrac16\tr(\Gamma+\Gamma\dg),\\
\Gamma &= \tr(\Sigma'B_1U)\,Q + \tr(\Sigma'B_2U)\,Q^2 + f_1\,U\Sigma' + f_2\,Q\,U\Sigma' + f_2\,U\Sigma'Q,\\
B_i &= b_{i0} + b_{i1}Q + b_{i2}Q^2,
\end{align}
with $b_{ij}$ derivatives of the $f_j$ by $c_0$, $c_1$ (in the code, \code{r01}..\code{r22} and the
\code{b}'s; $\xi_1(w)=(w\cos w-\sin w)/w^3$ by its series for $w<\tfrac14$). Then
\begin{align}
\Sigma &= \Sigma'\,e^{iQ} + C\dg\,i\Lambda + \sum(\text{terms of the staples that contain }U_\mu(x)),\\
i\Theta &= i\Lambda\,U .
\end{align}
The first kernel makes the first two terms and $i\Theta$; the second adds the six staple terms of
each plane, which read $i\Theta$ at the neighbours. On several GPUs $i\Theta$ is extended between
the two.

\subsection{\code{gauge\_force\_kernel} and \code{gauge\_traces\_kernel}}

\begin{align}
F_\mu(x) &= -\frac\beta3\,\mathrm{TA}\Big[U_\mu(x)\,\Sigma_\mu(x)\dg\Big],\\
\Sigma_\mu &= c_{\text{plaq}}\sum(\text{2 staples per plane}) + c_{\text{rect}}\sum(\text{6 rectangle staples per plane}),
\end{align}
each staple a product of links along a path (\code{walk}). The traces kernel sums per site
$\mathrm{Re}\,\tr$ of the 6 plaquettes and 12 rectangles; the action follows on the host after the
sum over sites.

\subsection{\code{links\_exp\_update\_kernel}: $U\leftarrow e^{\epsilon H}U$}

The Taylor series to order $n$ in nested (Horner) form:
\begin{align}
e^{\epsilon H}U &\approx \sum_{k=0}^n\frac{(\epsilon H)^k}{k!}U\\
&= U + \epsilon H\Big(U + \frac{\epsilon H}{2}\Big(U + \frac{\epsilon H}{3}\big(\cdots\big)\Big)\Big),
\end{align}
computed from the inside out:
\begin{equation}
u^{(0)} = U,\qquad u^{(k+1)} = U + \frac{\epsilon}{n-k}\,H\,u^{(k)},\qquad k=0,\dots,n-1 .
\end{equation}
Then reunitarised: $r_0\to r_0/|r_0|$, $r_1\to(r_1-(\bar r_0\cdot r_1)r_0)/|\cdot|$,
$r_2=\overline{r_0\times r_1}$.

\subsection{Other small kernels}

\code{links\_axpy} ($y\mathrel{+}=ax$), \code{links\_ta} ($s\,\mathrm{TA}(m)$), \code{links\_mul}
($ab$ or $a\dg b$), \code{links\_norm2} (per site $\sum$ reals$^2$, for $\tfrac12\sum\tr P^2$).

%======================================================================
\section{The multigrid: P}

\subsection{What P holds}

An aggregate keeps the whole fifth dimension and is split by the chirality of
$\Gamma_5=\gamma_5R$, $(R\psi)_s=\psi_{L_s-1-s}$. For a test vector $v$:
\begin{equation}
v^{(c)} = \tfrac12\big(1+(-1)^c\Gamma_5\big)v,\qquad c=0,1 .
\end{equation}
At fifth coordinate $s$:
\begin{equation}
v^{(c)}_s = \tfrac12\big(v_s + (-1)^c\gamma_5v_{L_s-1-s}\big).
\end{equation}
Now look at the mirror $L_s-1-s$:
\begin{align}
v^{(c)}_{L_s-1-s} &= \tfrac12\big(v_{L_s-1-s} + (-1)^c\gamma_5v_s\big)\\
&= (-1)^c\gamma_5\cdot\tfrac12\big((-1)^c\gamma_5v_{L_s-1-s} + v_s\big)\qquad(\gamma_5^2=1,\ ((-1)^c)^2=1)\\
&= (-1)^c\gamma_5\,v^{(c)}_s .
\end{align}
So the half $s<L_s/2$ determines the rest, and P is stored folded: $12L_s$ reals per (site, test
vector, chirality) instead of $24L_s$, in single or half precision. \code{fold\_fill\_kernel}
writes $\tfrac12(v_s\pm\gamma_5v_{L_s-1-s})$ for $s<L_s/2$.

\subsection{\code{fold\_restrict\_kernel}: coarse $=P\dg$fine}

For coarse site $X$, test vector $i$, chirality $c$:
\begin{align}
\text{coarse}(X,cn+i) &= \sum_{x\in X}\sum_{s=0}^{L_s-1} P(x,i,c)_s\dg\,f_s(x)\\
&= \sum_{x\in X}\sum_{s<L_s/2}\Big[P_s\dg f_s + P_{L_s-1-s}\dg f_{L_s-1-s}\Big]\\
&= \sum_{x\in X}\sum_{s<L_s/2}\Big[P_s\dg f_s + \big((-1)^c\gamma_5P_s\big)\dg f_{L_s-1-s}\Big]\\
&= \sum_{x\in X}\sum_{s<L_s/2}P_s\dg\Big[f_s + (-1)^c\gamma_5f_{L_s-1-s}\Big]\qquad(\gamma_5\dg=\gamma_5).
\end{align}
The last line is what the kernel computes: one read of the folded P per pair $(s,L_s-1-s)$. A thread
for each (coarse site, component), up to 8 vectors per read of P, which is most of the cost.

\subsection{\code{fold\_prolong\_kernel}: fine $=P$coarse}

\begin{align}
f_s(x) &= \sum_{c,i} P(x,i,c)_s\,\text{coarse}(X,cn+i),\\
f_{L_s-1-s}(x) &= \sum_{c,i}(-1)^c\gamma_5P(x,i,c)_s\,\text{coarse}(X,cn+i),
\end{align}
a thread for each (fine site, $s<L_s/2$, spin, colour), writing $s$ and its mirror.

\subsection{Orthonormalising P}

Per (aggregate, chirality), the $n$ vectors are made orthonormal over the aggregate. With
$W=[w_1\dots w_n]$:
\begin{equation}
G = W\dg W = LL\dg\ (\text{Cholesky}),\qquad Q = W\,(L\dg)^{-1}\ \Rightarrow\ Q\dg Q = L^{-1}GL^{-\dg} = 1 .
\end{equation}
\code{fold\_orthonormalise\_cholesky\_kernel} does this in double precision. Cholesky squares the
condition number, so where a pivot is below $10^{-12}$ of its diagonal the kernel gives up on that
site and \code{fold\_orthonormalise\_kernel} redoes it by modified Gram–Schmidt. The folded P holds
half of every inner product, so the result is scaled by $1/\sqrt2$.

%======================================================================
\section{The multigrid: the coarse operator}

\subsection{Probing: \code{hop\_split\_kernel}, \code{coarse\_units}, \code{coarse\_columns}}

With $e_j$ the unit vector of component $j$ on every coarse site:
\begin{align}
\phi_j &= P\,e_j,\\
A(X)_{\cdot j} &= \big[P\dg\,C\phi_j\big](X) + \sum_{\text{legs inside the block}}\big[P\dg\,\text{leg}\,\phi_j\big](X),\\
B_d(X)_{\cdot j} &= \big[P\dg\,\text{leg}_d\,\phi_j\big](X)\quad\text{on the sites whose leg }d\text{ leaves the block}.
\end{align}
\code{hop\_split\_kernel} computes one leg and sorts it: into \code{cross} where it comes from
another block, added to \code{inside} elsewhere. The restriction of \code{cross} is column $j$ of
$B_d$; \code{coarse\_columns} writes it straight into the matrices, in the precision they are kept
in. Several $j$ at a time, as the vectors of the fields.

\subsection{\code{coarse\_pass\_kernel}}

One kernel for every part of the coarse operator:
\begin{equation}
\text{out}(X) = M(X)\Big[\text{base}(X) + A(X)\,d(X) \pm \sum_{d}B_d(X)\,h\big(N[d][X]\big)\Big],
\end{equation}
each term optional. Choices give $D_c$, the Schur complement on the even sites
$S = A_{ee}-B_{eo}A_{oo}^{-1}B_{oe}$, its source and solution, the block inverses.
$D_c\dg=\Gamma_5D_c\Gamma_5$ ($\Gamma_5$: $-1$ on the second half of the components) is the same
kernel with a flag.

\textbf{Threads.} A block per (coarse site, vector); a warp per row; the lanes take pairs of
columns $(2\ell,2\ell+1),(2\ell+64,\dots)$, so a warp reads a row of a matrix as one stretch in
16-byte words. The fields at $X$ and its 8 neighbours are read once into shared memory. The sum of
a row is the lanes' parts added by a fixed butterfly, so the result does not depend on timing.

\subsection{\code{blocks\_inverse\_kernel}: Gauss–Jordan}

For each site, $[A\,|\,1]\to[1\,|\,A^{-1}]$, thread $r$ holding row $r$:
\begin{align}
&\text{for }k=0,\dots,\text{dim}-1:\\
&\quad p = \arg\max_{r\ge k}|a_{rk}|,\ \text{swap rows }k,p,\\
&\quad \text{row}_k \leftarrow \text{row}_k/a_{kk},\\
&\quad \text{row}_r \leftarrow \text{row}_r - a_{rk}\,\text{row}_k\quad(r\ne k).
\end{align}
A zero pivot sets a failure flag (a singular block).

\subsection{\code{blocks\_abs\_inverse\_gamma5\_kernel}: $|H|^{-1}$ by Jacobi}

$H$ = Hermitian part of $A\Gamma_5$. The cyclic Jacobi method: for each pair $(p,q)$,
\begin{align}
&\text{phase: } h_{pq} = |h_{pq}|e^{i\varphi}\ \Rightarrow\ \text{multiply column }q\text{ by }e^{-i\varphi},\ \text{row }q\text{ by }e^{i\varphi}\ (h_{pq}\text{ now real}),\\
&\text{rotation: } \tan2\vartheta = \frac{2h_{pq}}{h_{qq}-h_{pp}},\ \text{rotate rows and columns }p,q\ \Rightarrow\ h_{pq}=0,
\end{align}
repeated until the off-diagonal part vanishes, accumulating $V$. Then
\begin{equation}
|H|^{-1} = V\,|\Lambda|^{-1}V\dg,
\end{equation}
with eigenvalues below a cutoff times the largest left out (they come from dependent vectors).

\subsection{The exact adjoint: \code{coarse\_adjoint\_kernel}}

\begin{equation}
(A\dg u)(Y) = A(Y)\dg u(Y) + \sum_d B_d\big(N[\bar d][Y]\big)\dg\,u\big(N[\bar d][Y]\big)
\end{equation}
(derived in the other note); the kernel writes $B_d(X)\dg$ at the place of $Y$, entry by entry:
$(m\dg)_{ij} = \overline{m_{ji}}$.

%======================================================================
\section{The levels below the coarse lattice}

\begin{align}
\code{level\_basis}:&\quad Q_{\text{block}} = \text{Gram--Schmidt twice of the parts of the test vectors on a super-block and a half},\\
\code{level\_restrict}:&\quad y = Q\dg v\quad(\text{a warp per output, lanes along the sum}),\\
\code{level\_prolong}:&\quad v = Qy\ \text{or}\ v\mathrel{+}=Qy,\\
\code{level\_sources}:&\quad e_j = 1\ \text{at component }j\text{ of the sites of one colour},\\
\code{level\_probe}:&\quad A(t)_{\cdot j} = r_j(t)\ \text{if }t\text{ has the colour},\quad B_d(t)_{\cdot j} = r_j(t)\ \text{if }N[d][t]\text{ has it},\\
\code{level\_dense}:&\quad \text{the last level as one }n\times n\text{ matrix},\ n=\text{sites}\times\text{components},\\
\code{dense\_apply}:&\quad \text{out} = M\,r .
\end{align}
Why probing works: the sources of one colour are at least 3 apart in every direction, so
\begin{equation}
r_j(t) = \sum_{t'}D(t,t')e_j(t') = D(t,t'_*)\,\hat e_j\quad\text{with }t'_*\in\{t,\ t\pm\hat\mu\}\text{ the only source in reach}.
\end{equation}
The response at $t$ comes from one source only, and which one tells whether it is $A(t)$ or a
$B_d(t)$. The dense inverse (LU) and $|H|^{-1}$ of the last level are done by cuSOLVER
(\code{dense.cu}).

%======================================================================
\section{Linear algebra: \code{blas.cuh}}

\subsection{Fused operations}

A solver step usually needs several updates and a norm at once, e.g.\
\begin{equation}
r \leftarrow r - \alpha Ap,\qquad |r|^2 .
\end{equation}
\code{blas\_map\_kernel} and \code{blas\_reduce\_kernel} take the operation as a functor
(\code{Op}) and do it in one pass over the fields, every vector of a set at once.

\subsection{Deterministic reductions}

A sum over $N$ numbers is done in two stages, both in a fixed order:
\begin{align}
s_b &= \sum_{p\in\text{block }b}x_p\quad(\text{a tree inside each CUDA block, in double precision}),\\
S &= \sum_b s_b\quad(\code{blas\_finish\_kernel},\ \text{one block, a fixed tree}).
\end{align}
No atomic additions, so the result does not depend on which block finishes first: runs repeat bit
for bit. Only the final numbers go to the host.

\subsection{Sets of vectors}

\code{blas\_gram\_kernel}: all $\langle x_i,y_j\rangle$ in one pass. \code{blas\_rotate\_kernel}:
$x\leftarrow xC$ in place for a small $C$. \code{blas\_combine\_kernel}: $\text{out}_j=\sum_i x_ic_{ij}$
for larger sets, the outputs of a chunk in registers.

%======================================================================
\section*{Appendix: kernels removed}

Extracted from Bridge++ with the rest but never called by the library, and removed on 2026-10-24
with the C tests that exercised only them: the Schwarz domain decomposition (\code{dd.cuh}), an
older restriction, prolongation and coarse operator (\code{transfer.cuh}, \code{mdwf\_transfer.cuh},
\code{coarse.cuh}, \code{coarse\_eo.cuh}), multi-right-hand-side Wilson (\code{mrhs.cuh}), APE smearing
and a plaquette (\code{gauge.cuh}), the Dirac/chiral conversion, a plain real \code{axpy}, and an
unfinished Wilson face packing.

\end{document}
